
\section{\BookRecordIndex{T0247}模冒泡收敛}\label{sec:4-6}
\zhindex{momaopaoshoulian}{模冒泡收敛}

在第~\ref{sec:4-2} 节中，我们已经看到：当有限能量 \BookRecordIndex{T0003}$J$-全纯曲线列的导数趋于无穷时，
\BookRecordIndex{T0243}冒泡如何发生。本节还允许拉格朗日边界条件，并说明边界附近的\BookRecordIndex{T0243}冒泡如何产生边界值位于
\BookRecordIndex{T0219}拉格朗日子流形中的 $J$-全纯圆盘。由于 $J$-全纯圆盘或球面不可能具有任意小的能量
（命题~\ref{prop:4-1-4}），取适当子列后，\BookRecordIndex{T0243}冒泡只能发生在有限多个点附近。下一个定理
更精确地表述这种极限行为。

设 $(\Sigma,j_\Sigma,d\operatorname{vol}_\Sigma)$ 是一个复结构固定的黎曼曲面，它未必
紧；$(M,\omega)$ 是紧\BookRecordIndex{T0358}辛流形；$L\subset M$ 是紧\BookRecordIndex{T0219}拉格朗日子流形；
$J\in\mathcal J_\tau(M,\omega)$ 是被 $\omega$ \BookRecordIndex{T0372}驯服的\BookRecordIndex{T0088}殆复结构；而 $\hbar>0$ 是
命题~\ref{prop:4-1-4} 中的常数。

\begin{theorem}\label{thm:4-6-1}
设 $J^\nu\in\mathcal J_\tau(M,\omega)$ 是一列被 $\omega$ \BookRecordIndex{T0372}驯服的\BookRecordIndex{T0088}殆复结构，并在
$C^\infty$ 拓扑下收敛到 $J$；$\Omega^\nu$ 是一列递增且穷尽 $\Sigma$ 的开集；而
\[
  u^\nu:(\Omega^\nu,\Omega^\nu\cap\partial\Sigma)
  \longrightarrow(M,L)
\]
是一列 $J^\nu$-全纯曲线，满足
\[
  \sup_\nu E(u^\nu)<\infty.
\]
则存在一个子列（仍记作 $u^\nu$）、一条 \BookRecordIndex{T0003}$J$-全纯曲线
$u:(\Sigma,\partial\Sigma)\to(M,L)$，以及有限集
$Z=\{z_1,\ldots,z_\ell\}\subset\Sigma$，使得：
\begin{enumerate}[label=(\roman*),leftmargin=3em]
\item $u^\nu$ 在 $\Sigma\setminus Z$ 的紧子集上一致且连同所有阶导数一起收敛到
$u$；
\item 对每个 $j$ 以及每个满足
$B_\varepsilon(z_j)\cap Z=\{z_j\}$ 的 $\varepsilon>0$，极限
\[
  m_\varepsilon(z_j)
  :=\lim_{\nu\to\infty}E\bigl(u^\nu;B_\varepsilon(z_j)\bigr)
\]
\par\noindent 存在，且是 $\varepsilon$ 的连续函数，并且
\[
  m(z_j):=\lim_{\varepsilon\to0}m_\varepsilon(z_j)\ge\hbar;
\]
\item 对每个满足 $Z\subset\operatorname{int}(K)$ 的紧子集 $K\subset\Sigma$，
\[
  E(u;K)+\sum_{j=1}^\ell m(z_j)
  =\lim_{\nu\to\infty}E(u^\nu;K).
\]
\end{enumerate}
\end{theorem}

\begin{remark}\label{rem:4-6-2}
定理~\ref{thm:4-6-1} 中的收敛子列 $u^\nu$ 对每个点列
$z_j^\nu\to z_j$ 都满足
\[
  \lim_{\nu\to\infty}E\bigl(u^\nu;B_\varepsilon(z_j^\nu)\bigr)
  =m_\varepsilon(z_j).
\]
为看出这一点，固定任意小的 $\delta>0$，并注意当 $\nu$ 充分大时，
\[
  B_{\varepsilon-\delta}(z_j)
  \subset B_\varepsilon(z_j^\nu)
  \subset B_{\varepsilon+\delta}(z_j).
\]
所以对每个 $\delta>0$，该极限都落在区间
$[m_{\varepsilon-\delta}(z_j),m_{\varepsilon+\delta}(z_j)]$ 中；由函数
$\varepsilon\mapsto m_\varepsilon(z_j)$ 的连续性即得上述恒等式。
\end{remark}

\begin{remark}\label{rem:4-6-3}
在定理~\ref{thm:4-6-1} 中，$Z$ 的每一点都会发生\BookRecordIndex{T0243}冒泡。这意味着，对每个 $j$，都存在
一条非常值 $J$-全纯球面 $v:S^2\to M$，或一条非常值 $J$-全纯圆盘
$v:(B,\partial B)\to(M,L)$，使得对每个 $\varepsilon>0$，它的像都可以由
$u^\nu(B_\varepsilon(z_j))$ 中的点来逼近。
\end{remark}

定理~\ref{thm:4-6-1} 的证明依赖下面的引理~\ref{lem:4-6-5}。该引理断言：在\BookRecordIndex{T0251}能量密度
趋于无穷的每一点附近，都有至少为 $\hbar$ 的\BookRecordIndex{T0250}能量集中；而在每个这样的点，都会冒出一个
$J$-全纯球面或 $J$-全纯圆盘。这个引理的证明使用与第~\ref{sec:4-2} 节相同的\BookRecordIndex{T0148}共形重标度
论证。那里我们以\BookRecordIndex{T0251}能量密度的局部极大点为中心作重标度；但在当前情形中，定义域可能非紧，
因而也许不存在局部极大点。我们将看到，以“\BookRecordIndex{T0195}近似局部极大点”为中心作重标度便已足够。下一个
引理由 Hofer 给出，它说明如何找到这种点。

\begin{lemma}\label{lem:4-6-4}
设 $(X,d)$ 是度量空间，$f:X\to\mathbb R$ 是非负连续函数，并固定
$x\in X$ 和 $\delta>0$。假设闭球 $B_{2\delta}(x)$ 完备。则存在 $\zeta\in X$ 和
正常数 $\varepsilon\le\delta$，使得
\[
  \textup{(i)}\quad d(x,\zeta)<2\delta,
  \qquad
  \textup{(ii)}\quad \sup_{B_\varepsilon(\zeta)}f\le2f(\zeta),
  \qquad
  \textup{(iii)}\quad \varepsilon f(\zeta)\ge\delta f(x).
\]
\end{lemma}

\begin{proof}
假设结论不成立。我们断言，存在序列 $x_k\in X$，使得对 $k\ge0$，
\[
  x_0=x,
  \qquad
  d(x_k,x_{k+1})\le\frac\delta{2^k},
  \qquad
  f(x_{k+1})>2f(x_k),
\]
而这与完备性矛盾。用归纳法构造这个序列。先取 $x_0=x$。当
$\zeta=x_0$ 且 $\varepsilon=\delta$ 时，(i) 与 (iii) 成立。若 (ii) 不成立，则
\[
  \sup_{B_\delta(x_0)}f>2f(x_0),
\]
因而存在 $x_1\in X$，满足
\[
  d(x_0,x_1)\le\delta,
  \qquad f(x_1)>2f(x_0).
\]
现在令 $\zeta=x_1$、$\varepsilon=\delta/2$，重复这个论证。一般地，假设对某个
$k\ge1$，已经构造出适当的 $x_0,x_1,\ldots,x_k$。于是
\[
  d(x,x_k)<2\delta,
  \qquad
  f(x_k)>2^kf(x_0).
\]
所以当 $\zeta=x_k$ 且 $\varepsilon=\delta/2^k$ 时，(i) 与 (iii) 成立。若 (ii)
不成立，则
\[
  \sup_{B_{\delta/2^k}(x_k)}f>2f(x_k),
\]
因而存在 $x_{k+1}\in X$，满足
\[
  d(x_k,x_{k+1})\le\frac\delta{2^k},
  \qquad
  f(x_{k+1})>2f(x_k).
\]
这完成了归纳，也证明了引理~\ref{lem:4-6-4}。
\end{proof}

\begin{lemma}\label{lem:4-6-5}
设 $M,\omega,L,J^\nu,J$ 如定理~\ref{thm:4-6-1} 所述，并令
$\hbar=\hbar(M,\omega,L,J)$ 为命题~\ref{prop:4-1-4} 的常数。设
$\Omega\subset\mathbb H$ 是开集，且
\[
  u^\nu:(\Omega,\Omega\cap\mathbb R)\longrightarrow(M,L)
\]
是一列满足
\[
  \sup_\nu E(u^\nu;\Omega)<\infty
\]
的 $J^\nu$-全纯曲线。若存在点列 $z^\nu\to z_0\in\Omega$，使得
$|du^\nu(z^\nu)|\to\infty$，则对每个 $\varepsilon>0$，
\begin{equation}\label{eq:4-6-1}
  \liminf_{\nu\to\infty}
  E\bigl(u^\nu;B_\varepsilon(z_0)\cap\mathbb H\bigr)
  \ge\hbar,
\end{equation}
而且存在一条非常值 $J$-全纯球面 $v:S^2\to M$，或一条非常值 $J$-全纯圆盘
$v:(B,\partial B)\to(M,L)$。
\end{lemma}

\begin{proof}
取 $\varepsilon>0$，使得 $B_\varepsilon(z_0)\cap\mathbb H\subset\Omega$。把
引理~\ref{lem:4-6-4} 应用于函数
\[
  B_\varepsilon(z_0)\cap\mathbb H\longrightarrow\mathbb R,
  \qquad z\longmapsto|du^\nu(z)|,
\]
点 $x:=z^\nu$，以及常数
\[
  \delta:=\delta^\nu:=|du^\nu(z^\nu)|^{-1/2}.
\]
由此得到序列 $\zeta^\nu\in\mathbb C$ 和 $\varepsilon^\nu>0$，满足
\[
  \zeta^\nu\to z_0,
  \qquad
  \sup_{B_{\varepsilon^\nu}(\zeta^\nu)}|du^\nu|\le2c^\nu,
  \qquad
  \varepsilon^\nu\to0,
  \qquad
  \varepsilon^\nu c^\nu\to\infty,
\]
其中
\[
  c^\nu:=|du^\nu(\zeta^\nu)|.
\]
分成下列两种情形。（若 $\Omega\cap\mathbb R=\varnothing$，则总是情形 1。）

\medskip
\noindent\textbf{情形 1：}序列 $c^\nu\operatorname{Im}\zeta^\nu$ 无界；冒出一个
$J$-全纯球面。

不妨假设 $c^\nu\operatorname{Im}\zeta^\nu\to\infty$。极限点
$z_0=\lim_{\nu\to\infty}\zeta^\nu$ 可以在 $\Omega$ 的内部，也可以在
$\Omega\cap\mathbb R$ 中。不过，在后一种情形下，$\zeta^\nu$ 收敛到 $z_0$ 的速率
慢于导数趋于无穷的速率。由于
$c^\nu\operatorname{Im}\zeta^\nu\to\infty$，必要时可以缩小
$\varepsilon^\nu$，使得除了 $\varepsilon^\nu c^\nu\to\infty$ 外，还有
\[
  \varepsilon^\nu\le\operatorname{Im}\zeta^\nu,
  \qquad
  \varepsilon^\nu+|\zeta^\nu-z_0|\le\varepsilon.
\]
于是对每个 $\nu$，
\[
  B_{\varepsilon^\nu}(\zeta^\nu)
  \subset B_\varepsilon(z_0)\cap\mathbb H
  \subset\Omega.
\]
考虑\BookRecordIndex{T0411}重标度后的映射列
\[
  v^\nu(z):=u^\nu(\zeta^\nu+z/c^\nu).
\]
$v^\nu$ 定义在 $B_{\varepsilon^\nu c^\nu}$ 上，并且满足
\[
  \sup_{B_{\varepsilon^\nu c^\nu}}|dv^\nu|\le2,
  \qquad |dv^\nu(0)|=1,
\]
以及当 $R\le\varepsilon^\nu c^\nu$ 时，
\[
  E(v^\nu;B_R)
  =E\bigl(u^\nu;B_{R/c^\nu}(\zeta^\nu)\bigr).
\]
由于 $\varepsilon^\nu c^\nu\to\infty$，$v^\nu$ 在任意大的定义域上都有一致有界的
导数。由定理~\ref{thm:4-1-1}，存在一个子列（仍记作 $v^\nu$），它在
$\mathbb C$ 的紧子集上一致且连同所有阶导数一起收敛到一条 \BookRecordIndex{T0003}$J$-全纯曲线
$v:\mathbb C\to M$。极限映射 $v$ 具有有限能量；又因 $|dv(0)|=1$，它不是常值。
所以有限能量映射
\[
  \mathbb C\setminus\{0\}\longrightarrow M,
  \qquad z\longmapsto v(1/z),
\]
由定理~\ref{thm:4-1-2}(i) 可以光滑延拓过零点。于是 $v$ 延拓为黎曼球面
$S^2=\mathbb C\cup\{\infty\}$ 上的光滑 $J$-全纯映射，仍记作 $v$。由
命题~\ref{prop:4-1-4}，
\[
  E(v)\ge\hbar.
\]
所以，对每个实数 $\alpha\in(0,1)$，都存在常数 $R>0$，使
$E(v;B_R)>\alpha\hbar$。利用上面的区域包含，最后一项必须保留上半平面交集
\sourcefootnote{原书此处的能量链末项写作
$\liminf_\nu E(u^\nu;B_\varepsilon(z_0))$；但 $z_0$ 允许位于边界，原映射只定义在
$\mathbb H$ 中，而上面的包含关系与目标式~\eqref{eq:4-6-1} 均含
$B_\varepsilon(z_0)\cap\mathbb H$，故补上 $\cap\mathbb H$。}，从而
\[
  \begin{aligned}
  \alpha\hbar
  &<E(v;B_R)\\
  &=\lim_{\nu\to\infty}E(v^\nu;B_R)\\
  &=\lim_{\nu\to\infty}
    E\bigl(u^\nu;B_{R/c^\nu}(\zeta^\nu)\bigr)\\
  &\le\liminf_{\nu\to\infty}
    E\bigl(u^\nu;B_\varepsilon(z_0)\cap\mathbb H\bigr).
  \end{aligned}
\]
这对每个 $\alpha<1$ 都成立，因而\BookRecordIndex{T0326}推出式~\eqref{eq:4-6-1}。

\medskip
\noindent\textbf{情形 2：}序列 $c^\nu\operatorname{Im}\zeta^\nu$ 有界；冒出一个
$J$-全纯圆盘。

此时特别有 $\operatorname{Im}\zeta^\nu\to0$，所以
$z_0\in\Omega\cap\mathbb R$。取子列后，不妨假设
\[
  \varepsilon^\nu+|\operatorname{Re}\zeta^\nu-z_0|\le\varepsilon,
\]
而且极限
\[
  t_0:=\lim_{\nu\to\infty}c^\nu\operatorname{Im}\zeta^\nu\ge0
\]
存在。考虑\BookRecordIndex{T0411}重标度后的映射列
\[
  v^\nu(z):=u^\nu(\operatorname{Re}\zeta^\nu+z/c^\nu).
\]
由于
$B_{\varepsilon^\nu}(\operatorname{Re}\zeta^\nu)\cap\mathbb H\subset\Omega$，
$v^\nu$ 定义在 $B_{\varepsilon^\nu c^\nu}\cap\mathbb H$ 上。它满足
\[
  \sup_{B_{\varepsilon^\nu c^\nu}\cap\mathbb H}|dv^\nu|\le2,
  \qquad
  |dv^\nu(ic^\nu\operatorname{Im}\zeta^\nu)|=1,
\]
并且当 $R\le\varepsilon^\nu c^\nu$ 时，
\[
  E(v^\nu;B_R\cap\mathbb H)
  \le E\bigl(u^\nu;
    B_{R/c^\nu}(\operatorname{Re}\zeta^\nu)\cap\mathbb H\bigr).
\]
因此，$v^\nu$ 在任意大的定义域上都有一致有界的导数；再次由
定理~\ref{thm:4-1-1}，存在一个子列（仍记作 $v^\nu$），它在 $\mathbb H$ 的紧子集上
一致且连同所有阶导数一起收敛到一条 \BookRecordIndex{T0003}$J$-全纯曲线
\[
  v:(\mathbb H,\mathbb R)\longrightarrow(M,L).
\]
极限映射 $v:\mathbb H\to M$ 具有有限能量；又因 $|dv(it_0)|=1$，它不是常值。
因此，有限能量映射
\[
  \mathbb H\setminus\{0\}\longrightarrow M,
  \qquad z\longmapsto v(-1/z),
\]
由定理~\ref{thm:4-1-2}(ii) 可以光滑延拓过零点。于是定义在
$\mathbb D\setminus\{-1\}$ 上的函数
\[
  w(\zeta):=v\left(\frac{i(1-\zeta)}{1+\zeta}\right)
\]
延拓为一个 $J$-全纯圆盘
$w:(\mathbb D,\partial\mathbb D)\to(M,L)$。由命题~\ref{prop:4-1-4}，
\[
  E(v)=E(w)\ge\hbar.
\]
所以，对每个实数 $\alpha\in(0,1)$，都存在常数 $R>0$，使下述严格不等式成立
\sourcefootnote{原书此处写作
$E(v;B_R)>\alpha\hbar$；但本情形中的 $v$ 只定义在 $\mathbb H$ 上，且紧接着的能量链
也使用 $B_R\cap\mathbb H$，故据定义域唯一补上 $\cap\mathbb H$。}：
\[
  \begin{aligned}
  \alpha\hbar
  &<E(v;B_R\cap\mathbb H)\\
  &=\lim_{\nu\to\infty}E(v^\nu;B_R\cap\mathbb H)\\
  &\le\lim_{\nu\to\infty}
    E\bigl(u^\nu;B_{R/c^\nu}(\operatorname{Re}\zeta^\nu)
    \cap\mathbb H\bigr)\\
  &\le\liminf_{\nu\to\infty}
    E\bigl(u^\nu;B_\varepsilon(z_0)\cap\mathbb H\bigr).
  \end{aligned}
\]
这对每个 $\alpha<1$ 都成立，因而\BookRecordIndex{T0326}推出式~\eqref{eq:4-6-1}。引理~\ref{lem:4-6-5}
得证。
\end{proof}

\begin{proof}[定理~\ref{thm:4-6-1} 的证明]
若存在点列 $z^\nu\to z\in\Sigma$，使 $|du^\nu(z^\nu)|\to\infty$，就称
$z\in\Sigma$ 是映射列 $u^\nu$ 的\emph{奇异点}。由引理~\ref{lem:4-6-5}，每一列
满足 $\sup_\nu E(u^\nu)\le c$ 的 $J^\nu$-全纯曲线
$u^\nu:\Sigma\to M$ 至多有 $c/\hbar$ 个奇异点。

\emph{断言。}存在一个子列（仍记作 $u^\nu$），它的奇异点组成有限集
$Z=\{z_1,\ldots,z_\ell\}$，并且对每个紧集 $K\subset\Sigma\setminus Z$，
\[
  \sup_\nu\lVert du^\nu\rVert_{L^\infty(K)}<\infty.
\]

用归纳法证明这个断言。假设已经构造了一个子列（仍记作 $u^\nu$），其奇异点集包含
集合
\[
  Z_k=\{z_1,\ldots,z_k\}.
\]
（归纳从 $k=0$ 和 $Z_k=\varnothing$ 开始。）若对每个紧集
$K\subset\Sigma\setminus Z_k$，序列
$\lVert du^\nu\rVert_{L^\infty(K)}$ 都有界，则取 $\ell=k$，断言即得证明。否则，选取
一个紧集 $K\subset\Sigma\setminus Z_k$，使这个序列无界，并选择
$z_{k+1}^\nu\in K$，使得
\[
  \lVert du^\nu\rVert_{L^\infty(K)}
  =|du^\nu(z_{k+1}^\nu)|.
\]
必要时取子列，可设 $z_{k+1}^\nu$ 收敛到一点
$z_{k+1}\in\Sigma\setminus Z_k$，且
$|du^\nu(z_{k+1}^\nu)|\to\infty$。于是这个子列的奇异点集包含
\[
  Z_{k+1}:=Z_k\cup\{z_{k+1}\}.
\]
由于奇异点的可能数目已有上界，这个归纳过程必在有限步后终止。断言得证。

断言成立后，由定理~\ref{thm:4-1-1}，存在更进一步的子列（仍记作 $u^\nu$），它在
$\Sigma\setminus Z$ 的紧子集上一致且连同所有阶导数一起收敛到一条 \BookRecordIndex{T0003}$J$-全纯曲线
$u:\Sigma\setminus Z\to M$，并满足
$u(\partial\Sigma\setminus Z)\subset L$。\BookRecordIndex{T0214}可去奇点定理~\ref{thm:4-1-2} 断言，
$u$ 可以延拓为 \BookRecordIndex{T0003}$J$-全纯曲线
\[
  u:(\Sigma,\partial\Sigma)\longrightarrow(M,L).
\]

现在对 $j=1,\ldots,\ell$ 固定数 $\varepsilon_j>0$，使各球
$B_{\varepsilon_j}(z_j)$ 两两不交；再取一个子列，使极限
\[
  m_{\varepsilon_j}(z_j)
  :=\lim_{\nu\to\infty}
  E\bigl(u^\nu;B_{\varepsilon_j}(z_j)\bigr)
\]
对 $j=1,\ldots,\ell$ 都存在。这个子列满足定理的所有要求。特别地，由于 $u^\nu$
在\BookRecordIndex{T0170}环域 $B_{\varepsilon_j}(z_j)\setminus B_\varepsilon(z_j)$ 上收敛，(ii) 中的极限
对 $0<\varepsilon\le\varepsilon_j$ 都存在且连续。由命题~\ref{prop:4-1-4}，
\[
  m(z_j):=\lim_{\varepsilon\to0}m_\varepsilon(z_j)\ge\hbar.
\]

为证明 (iii)，固定 $\varepsilon\le\min_j\varepsilon_j$，并注意
\[
  \begin{aligned}
  E\left(u;K\setminus\bigcup_{j=1}^\ell B_\varepsilon(z_j)\right)
  &=\lim_{\nu\to\infty}E(u^\nu;K)
    -\sum_{j=1}^\ell\lim_{\nu\to\infty}
      E\bigl(u^\nu;B_\varepsilon(z_j)\bigr)\\
  &=\lim_{\nu\to\infty}E(u^\nu;K)
    -\sum_{j=1}^\ell m_\varepsilon(z_j).
  \end{aligned}
\]
令 $\varepsilon\to0$，得到
\[
  E(u;K)=\lim_{\nu\to\infty}E(u^\nu;K)
  -\sum_{j=1}^\ell m(z_j),
\]
正是所需结论。这就证明了定理~\ref{thm:4-6-1}。
\end{proof}

下面的引理处理一个将在第 5 章使用的小技术点。它说明：若曲线 $u$ 是一列 $J^\nu$-全纯
曲线 $u^\nu$ 的一致极限，则 $u^\nu$ 的所有阶导数也必一致收敛。注意，这里不需要再取
子列。

\begin{lemma}\label{lem:4-6-6}
设 $J^\nu,J$ 如定理~\ref{thm:4-6-1} 所述。设 $\Omega\subset\mathbb C$ 是开集，
$u^\nu:\Omega\to M$ 是一列 $J^\nu$-全纯曲线。再假设 $u:\Omega\to M$ 是连续函数，
且 $u^\nu$ 一致收敛到 $u$。则 $u$ 是 \BookRecordIndex{T0003}$J$-全纯曲线，而且 $u^\nu$ 在 $\Omega$ 的每个
紧子集上一致且连同所有阶导数一起收敛到 $u$。
\end{lemma}

\begin{proof}
先假设对每个紧子集 $K\subset\Omega$ 都有
\begin{equation}\label{eq:4-6-2}
  \sup_\nu\lVert du^\nu\rVert_{L^\infty(K)}<\infty.
\end{equation}
由标准\BookRecordIndex{T0330}椭圆自举方法，$u^\nu$ 的每个子列都有更进一步的子列，在 $\Omega$ 的紧子集上
一致且连同所有阶导数一起收敛到某条 \BookRecordIndex{T0003}$J$-全纯曲线。由假设，极限曲线总是 $u$。因此，
在这种情形下，结论来自通常的反证论证。更明确地说，若结论不成立，则存在一个子列
$\nu_j$ 和一个整数 $\ell$，使得对每个 $j$，
\[
  d_{C^\ell}(u^{\nu_j},u)\ge\delta
\]
对某个 $\delta>0$ 成立。但再取一个子列，便得到 $u^{\nu_j}$ 在 $C^\ell$ 拓扑下收敛
到 $u$，矛盾。这就证明了假设~\eqref{eq:4-6-2} 下的结论。

现在假设式~\eqref{eq:4-6-2} 不成立。引理~\ref{lem:4-6-5} 的证明表明，存在点列
$z^\nu\to z_0\in\Omega$ 和数列 $\varepsilon^\nu\to0$，使重标度
\[
  v^\nu(z)=u^\nu(z^\nu+\varepsilon^\nu z)
\]
有一个子列，一致且连同所有阶导数一起收敛到一条非常值 \BookRecordIndex{T0003}$J$-全纯曲线
$v:B_1\to M$。另一方面，由 $u^\nu$ 一致收敛到 $u$，对任意 $\varepsilon>0$，都存在
$\nu_0$，使得当 $\nu\ge\nu_0$ 时，
\[
  u^\nu(B_{\varepsilon^\nu}(z^\nu))
  \subset B_\varepsilon(u(z_0)).
\]
所以极限曲线 $v:B_1\to M$ 的像对每个 $\varepsilon>0$ 都包含在球
$B_\varepsilon(u(z_0))$ 中，从而 $v\equiv u(z_0)$。这与 $v$ 非常值矛盾。
引理~\ref{lem:4-6-6} 得证。
\end{proof}

定理~\ref{thm:4-6-1} 还相当粗略：它只断言\BookRecordIndex{T0275}气泡存在，却没有完整描述所有冒出的全纯
球面或圆盘，也没有说明映射列 $u^\nu$ 如何收敛到这一族\BookRecordIndex{T0275}气泡。为了更好地理解这种现象，
必须细化上述重标度论证。下一节将在 $J$-全纯球面的情形中进行这一讨论。
